% ============================================================
% CHAPTER 2 - LITERATURE REVIEW
% ============================================================
\chapter{Literature Review}
\label{ch:literature}

\section{Introduction}

This chapter surveys the work on which the present system is built. It proceeds
from the clinical characterisation of dysarthric speech
(Section~\ref{sec:lit-dysarthria}), through the corpora available for research
(Section~\ref{sec:lit-corpora}), to the two technical traditions that intersect
in this project: automatic pronunciation assessment
(Section~\ref{sec:lit-mdd}) and dysarthric speech processing
(Section~\ref{sec:lit-dysarthric-asr}). It then reviews the enabling
components used in the pipeline --- grapheme-to-phoneme conversion
(Section~\ref{sec:lit-g2p}), forced alignment
(Section~\ref{sec:lit-alignment}), self-supervised representations
(Section~\ref{sec:lit-ssl}) and the acoustic features and neural architectures
employed for phoneme-level classification
(Sections~\ref{sec:lit-features}--\ref{sec:lit-architectures}). It closes with
a discussion of labelling strategies for weakly supervised speech assessment
(Section~\ref{sec:lit-labelling}), a comparative summary
(Section~\ref{sec:lit-summary}) and a statement of the research gap
(Section~\ref{sec:lit-gap}).

\section{Dysarthric Speech: Clinical and Acoustic Characterisation}
\label{sec:lit-dysarthria}

Darley, Aronson and Brown's foundational perceptual studies established the
classification of dysarthria into flaccid, spastic, ataxic, hypokinetic,
hyperkinetic and mixed types, each associated with a distinct lesion site and
a distinct cluster of perceptual deviations \cite{darley1969differential}.
Duffy's subsequent synthesis \cite{duffy2019motor} remains the standard
clinical reference and frames dysarthria as a disorder of the execution of
speech movements, distinguishing it from apraxia of speech (a disorder of motor
planning) and from the language-level impairments of aphasia.

From a signal-processing perspective, four families of acoustic consequence are
consistently reported.

\textbf{Vowel space compression.} Kim, Hasegawa-Johnson and Perlman
\cite{kim2011vowel} showed on UASpeech data that the area of the vowel
quadrilateral in the $F_1$--$F_2$ plane shrinks with dysarthria severity, and
that this compression correlates with intelligibility. Acoustically distinct
vowels migrate towards a central, schwa-like target, which is why reduced
vowels are difficult to discriminate in impaired speech.

\textbf{Imprecise consonant articulation.} Stop consonants are released weakly
or not at all, voice-onset time becomes irregular, and fricatives lose the
high-frequency energy that carries their identity. Consonant imprecision is
routinely the single strongest perceptual correlate of reduced intelligibility
\cite{ansel1992acoustic}.

\textbf{Prosodic and temporal disruption.} Speaking rate is reduced, segment
durations lengthen and become more variable, and inter-word pauses increase
\cite{tjaden2003characteristics}. Segment duration in dysarthric speech is
therefore an unreliable cue and a poorly behaved feature dimension.

\textbf{High within-speaker variability.} Repetitions of the same word by the
same dysarthric speaker differ far more than repetitions by a control speaker.
This is a defining property of the population and it has direct methodological
consequences: any labelling scheme that assumes a stable acoustic target per
speaker will be noisy, and any model trained on such data must tolerate high
intra-class variance.

The final point deserves emphasis, because it distinguishes dysarthric errors
from the errors studied in the mainstream pronunciation-assessment literature.
A second-language learner who cannot produce \ph{θ} substitutes a
\emph{different but well-formed} phoneme. A dysarthric speaker attempting
\ph{θ} typically produces a token that is a degraded version of \ph{θ} --- not
a clean substitute. The distinction is between categorical substitution and
gradient distortion, and it implies that a classifier over phoneme identity is
the wrong formulation for this population; what is required is a measure of
\emph{deviation from a clear realisation of the intended target}.

\section{Corpora for Dysarthric Speech Research}
\label{sec:lit-corpora}

Two corpora dominate English-language dysarthria research.

\textbf{UASpeech} \cite{kim2008dysarthric}, released by the University of
Illinois Urbana--Champaign, contains isolated-word recordings from 15 speakers
with cerebral-palsy-related dysarthria spanning a wide intelligibility range,
together with 13 age-matched control speakers. The vocabulary comprises digits,
computer commands, radio alphabet letters, common words and uncommon words, and
each item is recorded in multiple blocks through a seven-microphone array. Its
decisive property for the present work is that it is \emph{paired}: for a given
word, both dysarthric and healthy control productions exist. This makes it
possible to define ``correct'' empirically, by reference to a healthy
production of the same word, rather than by an arbitrary statistical criterion.

\textbf{TORGO} \cite{rudzicz2012torgo} pairs acoustic recordings with
electromagnetic articulography, giving direct measurement of articulator
positions for speakers with cerebral palsy and amyotrophic lateral sclerosis.
It is smaller than UASpeech and its value lies chiefly in articulatory-inversion
research; it was not adopted here because the present method requires paired
matched-word control productions at scale.

Neither corpus provides clinician-annotated \emph{per-phoneme error labels}.
This absence is the single most important constraint on supervised
phoneme-level work in this domain and it is the reason the labelling strategy,
rather than the classifier, is the substantive methodological problem addressed
in this dissertation.

\section{Mispronunciation Detection and Diagnosis}
\label{sec:lit-mdd}

\subsection{Goodness of Pronunciation and its descendants}

The Goodness of Pronunciation (GOP) score introduced by Witt and Young
\cite{witt2000phone} remains the reference formulation for automatic
pronunciation assessment. For a phone segment $o$ aligned to canonical phone
$p$, GOP is the duration-normalised log posterior
\begin{equation}
\mathrm{GOP}(p) = \frac{1}{|o|}\,
\log \frac{P(o \mid p)}{\max_{q \in Q} P(o \mid q)},
\label{eq:gop}
\end{equation}
where $Q$ is the full phone inventory. A segment whose acoustics are better
explained by some competing phone than by the canonical one receives a low
score and is flagged. GOP is elegant, interpretable and requires no
mispronunciation-labelled training data, and it has been extended repeatedly:
with discriminatively trained acoustic models, with DNN posteriors in place of
HMM--GMM likelihoods \cite{hu2015improved}, and with transfer from
self-supervised encoders.

Its assumption, however, is precisely the one that fails for dysarthria.
Equation~\eqref{eq:gop} measures whether some \emph{other phoneme} explains the
segment better. A gradient distortion is explained poorly by \emph{every}
phoneme in the inventory, so both numerator and denominator are small and the
ratio is uninformative. GOP detects substitution well and distortion badly.

\subsection{End-to-end and neural MDD}

Neural approaches reformulate MDD as sequence labelling. Leung et al.
\cite{leung2019cnn} applied a CNN--RNN--CTC model to non-native English,
recognising the actually-produced phone sequence and comparing it to the
canonical sequence, and reported improvements over GOP baselines. Li et al.
\cite{li2016mispronunciation} combined acoustic--phonemic and
acoustic--graphemic models to improve diagnostic feedback. More recent work
fine-tunes \texttt{wav2vec\,2.0} for phone recognition and derives error
verdicts from the alignment between recognised and canonical sequences
\cite{xu2021explore}, and transformer-based MDD systems now dominate the
benchmark leaderboards on L2-ARCTIC.

Two limitations of this line of work matter here. It is trained and evaluated
almost exclusively on non-native learner speech, whose error distribution
differs from that of motor speech disorders in the manner described in
Section~\ref{sec:lit-dysarthria}; and it requires phone-level annotation of
what was \emph{actually} produced, which for dysarthric speech is both
expensive and unreliable, because trained annotators disagree substantially
when asked to transcribe distorted tokens.

\section{Automatic Recognition and Assessment of Dysarthric Speech}
\label{sec:lit-dysarthric-asr}

Research on dysarthric speech has concentrated on two tasks.

\textbf{Recognition.} The goal is to transcribe dysarthric speech accurately
despite acoustic mismatch with models trained on typical speech. Speaker
adaptation, data augmentation by tempo and formant perturbation, and
severity-aware model selection have all been shown to reduce word error rate
\cite{shor2019personalizing,geng2022speaker}. Google's Project Euphonia
demonstrated that personalised models trained on a modest amount of a
particular speaker's data can dramatically outperform speaker-independent
systems \cite{macdonald2021disordered}. This literature optimises transcription
accuracy, not diagnostic feedback.

\textbf{Severity and intelligibility estimation.} A second line predicts an
utterance-level or speaker-level intelligibility score, typically by regression
over utterance embeddings \cite{janbakhshi2019pathological}. These systems
answer ``how impaired is this speaker'' rather than ``which sound in this word
failed''. The output is a scalar, and scalars do not tell a person practising
at home what to practise.

The gap between these two lines is the space this dissertation occupies:
retaining the known target word (as in assessment) and producing a per-phoneme
verdict (as in MDD), for a population whose errors are gradient distortions
rather than categorical substitutions.

\section{Grapheme-to-Phoneme Conversion}
\label{sec:lit-g2p}

Deriving the canonical phoneme sequence for a target word is a prerequisite for
alignment. Approaches range from pronunciation lexicons such as CMUdict, which
are exact but closed-vocabulary, through joint-sequence n-gram models
\cite{bisani2008joint}, to neural sequence-to-sequence models
\cite{yao2015sequence}.

This work uses \texttt{espeak-ng}, an open-source formant-synthesis engine whose
rule-based front end emits IPA transcriptions for arbitrary input text. Three
properties motivated the choice: it is open-vocabulary, so a user may practise
any word including names and nonce words; its IPA inventory for the
\texttt{en-us} voice was verified to coincide with the vocabulary of the chosen
alignment model, so no symbol mapping table is required; and it is a small
native binary that adds negligible latency and memory footprint to a deployed
service. Its limitation is that it is rule-based and therefore mispronounces
some irregular English orthography, and that it emits a single canonical
pronunciation without alternatives.

\section{Forced Alignment}
\label{sec:lit-alignment}

Forced alignment determines the time boundaries of a known symbol sequence
within an utterance. The classical solution is Viterbi decoding through an
HMM constrained to the target sequence, as implemented in the Montreal Forced
Aligner \cite{mcauliffe2017montreal} and in HTK-based toolkits. These systems
are accurate on typical speech but depend on acoustic models trained on typical
speech, and their accuracy degrades on disordered speech.

The modern alternative aligns using the frame-wise posteriors of a CTC-trained
neural phone recogniser. Given the log-posterior matrix
$\log P(k \mid t)$ over symbols $k$ and frames $t$, and a target token sequence
$\mathbf{y}$, the CTC forced-alignment algorithm finds
\begin{equation}
\boldsymbol{\pi}^{*} = \arg\max_{\boldsymbol{\pi} \in \mathcal{B}^{-1}(\mathbf{y})}
\sum_{t=1}^{T} \log P(\pi_t \mid t),
\label{eq:ctc-align}
\end{equation}
where $\mathcal{B}$ is the CTC collapsing function that removes blanks and
repeated labels, so that $\mathcal{B}^{-1}(\mathbf{y})$ is the set of all frame
paths consistent with $\mathbf{y}$. The maximisation is solved exactly by
dynamic programming in $O(T|\mathbf{y}|)$ time. \texttt{torchaudio} provides an
optimised implementation \cite{yang2022torchaudio}, and it is this route that
the present work adopts.

The critical property of forced alignment for this work is that it is
\emph{constrained}: it is told what the target sequence is and must place it in
time. It therefore continues to produce boundaries even when the speech is so
degraded that free recognition would fail --- which is exactly the operating
regime of dysarthric speech. Its complementary weakness is that it never
declines to align: a phoneme that was entirely omitted by the speaker will
still be assigned some (short, low-confidence) interval. The alignment
confidence score returned alongside each boundary is retained in this work as
a diagnostic of that condition.

\section{Self-Supervised Speech Representations}
\label{sec:lit-ssl}

\texttt{wav2vec\,2.0} \cite{baevski2020wav2vec} learns speech representations by
masking spans of a latent feature sequence and solving a contrastive task over
quantised targets, using unlabelled audio only. Fine-tuned with a CTC head on a
small amount of labelled data, it approaches or exceeds fully supervised
systems. HuBERT \cite{hsu2021hubert} and WavLM \cite{chen2022wavlm} refine the
pre-training objective; Whisper \cite{radford2023robust} takes the alternative
route of large-scale weakly supervised multilingual training and is notable for
robustness to accent, noise and domain shift.

Two such models are used here, in distinct roles. The multilingual phone
recogniser \texttt{wav2vec2-lv-60-espeak-cv-ft} \cite{xu2022simple}, trained to
output IPA phone symbols, supplies the emission matrix for the forced alignment
of Equation~\eqref{eq:ctc-align}. Separately, \texttt{whisper-tiny.en} is used
in the deployed service purely as a target-word verifier --- to distinguish
``the user said a completely different word'' from ``the user mispronounced the
target word'' --- a role in which its known tendency to normalise disordered
speech towards plausible words is acceptable because only a coarse similarity
decision is required.

\section{Acoustic Features for Phoneme-Level Classification}
\label{sec:lit-features}

Mel-frequency cepstral coefficients (MFCCs) were for decades the default
representation for speech classification, being compact and largely
decorrelated. For convolutional models, however, the log-mel spectrogram is now
generally preferred: the discrete cosine transform that produces MFCCs destroys
the local time--frequency structure that convolutional kernels are designed to
exploit \cite{purwins2019deep}.

The log-mel spectrogram is obtained by applying a mel filterbank
$\{H_m\}_{m=1}^{M}$ to the short-time power spectrum and taking logarithms,
\begin{equation}
S(m,t) = 10\log_{10}\!\left(\sum_{k} H_m(k)\,\bigl|X(k,t)\bigr|^{2}\right),
\label{eq:logmel}
\end{equation}
where the mel scale
$\mathrm{mel}(f) = 2595\log_{10}\!\left(1 + f/700\right)$
compresses frequency in accordance with human auditory resolution. Log-mel
patches are the representation adopted in this work, at 64 mel bands and a
fixed 32-frame width per phoneme.

Because the log-mel representation is in decibels, its numerical scale depends
on the recording level. Global $z$-score normalisation using statistics
computed over the training corpus is therefore applied, and --- as
Chapter~\ref{ch:implementation} discusses --- reproducing those exact
statistics at inference time is a hard requirement for the deployed system,
not an optional refinement.

\section{Neural Architectures for Speech Classification}
\label{sec:lit-architectures}

\subsection{Convolutional networks}

Two-dimensional convolution over a time--frequency image exploits the locality
and translation equivariance of spectro-temporal patterns: formant transitions,
fricative noise bands and stop bursts are local structures whose identity is
largely independent of absolute position. Abdel-Hamid et al.
\cite{abdelhamid2014convolutional} established the value of convolutional
layers for speech recognition, and CNNs over log-mel inputs remain the standard
front end for audio classification \cite{hershey2017cnn}.

\subsection{Recurrent networks}

A phoneme is not a static pattern; its identity and its quality are carried by
\emph{transitions}. Long Short-Term Memory (LSTM) networks
\cite{hochreiter1997long} model such temporal dependence through gated memory,
and Gated Recurrent Units (GRUs) \cite{cho2014learning} achieve comparable
modelling capacity with fewer parameters. Bidirectional variants
\cite{schuster1997bidirectional,graves2005framewise} process the sequence in
both directions, so that the representation at each step is conditioned on both
preceding and following context. For a phoneme patch that has already been
extracted in full, bidirectionality is unambiguously appropriate: the entire
segment is available at inference time, and the release of a stop is only
interpretable in the light of its closure.

\subsection{Hybrid CNN--RNN models}

The combination --- convolutional layers extracting local spectro-temporal
features, followed by a recurrent layer integrating them along the time axis
--- is the dominant design for audio classification tasks of this shape
\cite{choi2017convolutional}, and variants of it have been reported for
dysarthric speech classification \cite{joshy2022automated}. This is the family
from which the five candidates compared in Chapter~\ref{ch:results} are drawn.

\subsection{Transformer alternatives}

Transformer encoders now dominate large-scale speech modelling. They were not
adopted for the classifier here for three reasons: the input is a very short
fixed-size patch (16 time steps after pooling), a regime in which self-attention
offers little advantage over recurrence; the labels are proxy labels derived
from acoustic distance, so the effective supervision is noisy and a
high-capacity model would fit that noise; and the deployment target is free-tier
CPU hardware, where a 0.4-million-parameter recurrent model is preferable to a
transformer by an order of magnitude in latency. The pre-trained transformer
capacity that the system \emph{does} use is confined to the alignment stage,
where it is applied as a frozen feature extractor.

\section{Labelling Strategies for Weakly Supervised Speech Assessment}
\label{sec:lit-labelling}

Where per-token annotation is unavailable, labels must be derived. Three
families of strategy appear in the literature.

\textbf{Reference-comparison methods} compare a production against a reference
realisation of the same content, typically with dynamic time warping over
spectral features. This is the classical basis of pronunciation scoring in
computer-aided language learning \cite{neri2002feedback}, and it is the family
to which the present method belongs.

\textbf{Model-posterior methods} such as GOP (Equation~\eqref{eq:gop}) derive
labels or scores from an acoustic model's own confidence.

\textbf{Distribution-thresholding methods} compute a continuous deviation
statistic for every token and declare a fixed proportion of the distribution
erroneous. This is attractive because it needs no reference and no acoustic
model, and it guarantees a workable class balance.

It is the third family that contains the failure mode this dissertation
documents. If a percentile cutoff is applied to a distribution pooled over both
impaired and healthy speakers, then by definition a fixed fraction of the
healthy tokens is assigned to the error class, however clear those tokens
actually were. The label noise so introduced is not random: it is systematic
and concentrated exactly on the examples that should define correctness. A
model trained on such labels learns the base rate, not the phenomenon. The
control-baseline strategy of Section~\ref{sec:method2} is a reference-comparison
method designed specifically to eliminate this pathology, by restricting the
percentile estimation to the impaired sub-population and treating the healthy
sub-population as a set of guaranteed-correct anchors.

\section{Comparative Summary of Related Work}
\label{sec:lit-summary}

Table~\ref{tab:related-work} positions the present work against the closest
existing systems.

\begin{table}[H]
\centering
\caption[Comparison of related work]{Comparison of representative related work
against the present system.}
\label{tab:related-work}
\small
\begin{tabularx}{\textwidth}{@{}p{3.1cm}p{2.5cm}p{2.4cm}X@{}}
\toprule
\textbf{Work} & \textbf{Population} & \textbf{Output} & \textbf{Limitation for the present problem} \\
\midrule
Witt \& Young (GOP) \cite{witt2000phone} & L2 learners & Per-phone score & Assumes competing phone explains segment; blind to gradient distortion \\
Li et al. \cite{li2016mispronunciation} & L2 learners & Phone-level diagnosis & Requires annotated produced-phone sequences \\
Leung et al. \cite{leung2019cnn} & L2 learners & Recognised phone sequence & Free phone recognition unreliable on severely disordered speech \\
Kim et al. \cite{kim2011vowel} & Dysarthric & Vowel-space metrics & Analysis only; vowels only; not a per-phoneme detector \\
Project Euphonia \cite{macdonald2021disordered} & Disordered & Transcript & Optimises WER; no per-phoneme feedback \\
Janbakhshi et al. \cite{janbakhshi2019pathological} & Pathological & Intelligibility score & Utterance-level scalar; not actionable per sound \\
\textbf{This work (\systemName)} & \textbf{Dysarthric / unclear} & \textbf{Per-phoneme verdict + probability} & \textbf{Proxy labels from acoustic distance; not clinician-validated} \\
\bottomrule
\end{tabularx}
\end{table}

\section{Research Gap}
\label{sec:lit-gap}

The review supports four observations.

\begin{enumerate}[leftmargin=*]
  \item Automatic pronunciation assessment is mature but is formulated around
        categorical substitution errors in second-language learners, and its
        central scoring criterion (GOP) is structurally insensitive to the
        gradient distortions that characterise dysarthria.
  \item Dysarthric speech processing is mature but is directed at transcription
        accuracy or at scalar severity estimation, neither of which produces
        the per-sound feedback that unsupervised therapy practice requires.
  \item No public corpus provides clinician-annotated per-phoneme error labels
        for dysarthric speech, so supervised phoneme-level detection depends on
        derived labels; yet the derivation strategy is rarely scrutinised, and
        the most convenient strategy --- global percentile thresholding ---
        systematically mislabels healthy speech.
  \item Even where research prototypes exist, they are seldom carried through
        to a deployable artefact that a person could actually use for practice
        on a phone or laptop.
\end{enumerate}

This dissertation addresses all four: it targets dysarthric rather than L2
speech, produces per-phoneme rather than scalar output, treats the derivation
of the ground truth as the primary methodological problem and solves it with a
control-baseline strategy, and delivers the result as a deployed application.
